\honest{Collective Apathy and the Internet}{Eliezer Yudkowsky}{2009}

Yesterday I covered the bystander effect: a group of bystanders is less likely to act than one bystander. The standard explanations are pluralistic ignorance and diffusion of responsibility.

This ``may be'' a sign that our hunter-gatherer instincts for coordinating, built for groups of people we knew, fail in modern conditions. When all the subjects know each other, the bystander effect diminishes. \nb{Yesterday the author said this with ``If I recall correctly.'' The source bears it out.}

With mock astonishment at my own daring, I observe that people ``seem to have a hard time reacting constructively to problems encountered over the Internet.'' I give no example. Perhaps our instincts are not tuned for groups of strangers of unknown size, with no eye contact, no real-time talk, no ties of mutual help, the cover of anonymity, a crowd of other people doing nothing, and no voiced plea for help. \nb{This had been tested. In 2000 Patrick Markey found that in chat groups help took longer the more people were present, and that the effect was ``virtually eliminated'' when the request named one person.}

I have no brilliant solution. Online activism tools ``tend to be'' of the ``sign this petition!'' kind. \nb{PledgeBank, running since 2005, was built to get people past ``the fear of acting alone'': you pledged to act if enough others would.} Some ideas: a video of someone asking for help; names and photos of the first helpers; video thank-yous; referral codes that show how many people you brought in; bounties for subtasks that can be checked online. The same applies to open-source code and even petitions, though ``money is usually the most effective.'' \nb{Donations and code are contributions to a shared good, the case Mancur Olson analyzed in 1965. Rewards that go only to those who contribute, like these, are what Olson called selective incentives; Olson argued that large groups need them in order to act at all.}

Mostly I hand you an open problem: help groups of strangers form effective task forces online. All of Wikipedia, by Clay Shirky's estimate, is 1/2,000 of the time Americans spend watching television. \nb{Shirky's comparison is with a single year of television.} There is plenty of fuel, if only there were an engine.
